\section{Simulazione a eventi discreti per reti veicolari}
\label{env:discrete_event}

La simulazione a eventi discreti (DES, \emph{Discrete Event Simulation}) è il paradigma più utilizzato per lo studio di protocolli di rete e sistemi distribuiti complessi, poiché consente di rappresentare l’evoluzione del sistema come sequenza ordinata di eventi temporizzati. In tale paradigma, il tempo simulato avanza da un evento al successivo (trasmissione, ricezione, timeout, aggiornamento di stato, ecc.), evitando un aggiornamento uniforme a passo fisso dell’intero sistema quando non necessario \cite{varga2008overview}. Questo approccio è particolarmente adatto ai sistemi V2X, in cui coesistono eventi di rete molto frequenti e eventi di mobilità/controllo su scale temporali differenti.

Nel caso del platooning cooperativo, la DES è utile perché permette di modellare in modo esplicito l’interazione tra:
\begin{itemize}
\item la generazione periodica di messaggi di controllo e beacon;
\item la contesa per il canale e le perdite di pacchetto;
\item i timeout e i meccanismi di \emph{fallback}/\emph{recovery};
\item l’aggiornamento della dinamica del veicolo e del controllore.
\end{itemize}
L’aspetto cruciale non è soltanto simulare bene la rete oppure la mobilità, ma preservare la coerenza temporale tra i due domini, così da misurare l’effetto delle condizioni di comunicazione sul comportamento del platoon.

Rispetto a una modellazione puramente analitica o a simulazioni separate (rete da una parte, traffico dall’altra), una DES accoppiata consente di osservare effetti emergenti difficilmente catturabili con modelli semplificati: ad esempio degradazioni della stabilità dovute a burst di perdite, impatti degli handover sulla continuità della comunicazione, oppure retroazioni tra aumento della distanza di sicurezza e qualità del link disponibile. Questo è precisamente il tipo di fenomeni che motivano l’uso di toolchain integrate nelle ricerche su cooperative driving e platooning \cite{sommer2011bidirectionally,segata2014plexe,segata2023multi-technology}.

In questa tesi, la simulazione a eventi discreti costituisce quindi il livello di orchestrazione su cui vengono innestati i modelli di rete (OMNeT++ con relative estensioni), la mobilità stradale microscopica (SUMO) e la logica di platooning.
